% *** ก่อนส่ง: ลบข้อความตัวอย่าง (\ldots) ทั้งหมด แล้วเขียนเนื้อหาของตนเอง ***
% =============================================================
% บทที่ 5 - บทสรุป
% แก้ไขเนื้อหาในบทนี้ได้ที่ไฟล์ chapter5.tex
% =============================================================

% -----------------------------------------------------------
\section{สรุปผลการดำเนินโครงงาน}
% -----------------------------------------------------------

โครงงานนี้ได้พัฒนาระบบ Pacegasus มีวัตถุประสงค์เพื่อออกแบบและพัฒนาแอปพลิเคชันวิ่งที่ผสาน Adaptive Training Engine เข้ากับระบบ Gamification และระบบสังคมบนพื้นฐานทฤษฎีการกำหนดใจตนเอง (SDT) ผลการดำเนินงานสรุปเทียบกับวัตถุประสงค์ได้ดังตารางที่ \ref{tab:obj_result}

\begin{table}[H]
\centering
\caption{สรุปผลเทียบกับวัตถุประสงค์}
\label{tab:obj_result}
\small
\begin{tabular}{|c|p{4.3cm}|p{6.8cm}|p{1.7cm}|}
\hline
\textbf{ข้อ} & \textbf{วัตถุประสงค์} & \textbf{ผลที่ได้} & \textbf{สถานะ} \\
\hline
1 & พัฒนาแอปพลิเคชันมือถือสำหรับเดินและวิ่งที่ติดตามและบันทึกข้อมูลการเคลื่อนไหวของผู้ใช้ & พัฒนาต้นแบบ Pacegasus ครอบคลุมการสมัครสมาชิก การตั้งเป้าหมายและแผนฝึก การเริ่มและบันทึกการวิ่ง การแสดงผล และการประเมินหลังวิ่ง & \textbf{บรรลุบางส่วน} \\ \hline
2 & พัฒนาเอนจินปรับแผนฝึกอัตโนมัติจากระยะทาง เพซ และความสม่ำเสมอของผู้ใช้ & รองรับแผนฝึก 3 ระดับ พร้อมกฎ Weekly Cap และการปลดล็อกระดับ ส่วนการปรับแผนอัตโนมัติด้วย Session-RPE, ACWR และ Machine Learning เป็นแนวทางต่อยอด & \textbf{บรรลุบางส่วน} \\ \hline
3 & เพิ่มแรงจูงใจในการออกกำลังกายด้วยเกมมิฟิเคชันและการมีปฏิสัมพันธ์ทางสังคม & พัฒนาระบบภารกิจ รางวัล และความก้าวหน้า พร้อมออกแบบโครงสร้าง Club, Friends และ Leaderboard ซึ่งบางส่วนยังเป็นแนวทางต่อยอด & \textbf{บรรลุบางส่วน} \\
\hline
\end{tabular}
\end{table}

จากตารางที่ \ref{tab:obj_result} โครงงานนี้บรรลุวัตถุประสงค์ทั้ง 3 ข้อในระดับบางส่วน โดยสามารถพัฒนาฟังก์ชันหลักของระบบได้ตามขอบเขตที่กำหนด และยังมีฟังก์ชันบางส่วนที่สามารถพัฒนาต่อยอดได้ในอนาคต

% -----------------------------------------------------------
\section{ข้อจำกัดของระบบ}
% -----------------------------------------------------------

\subsection{ข้อจำกัดด้านการทำงาน}
\subsubsection{ความแม่นยำของข้อมูลการวิ่งขึ้นอยู่กับคุณภาพสัญญาณ GPS และเซนเซอร์ของอุปกรณ์}
\subsubsection{การปรับแผนการฝึกอาศัยข้อมูลที่ผู้ใช้บันทึกและประเมินตนเอง จึงอาจแตกต่างตามความถูกต้องของข้อมูลที่ผู้ใช้ให้}

\subsection{ข้อจำกัดด้านฮาร์ดแวร์และซอฟต์แวร์}
\subsubsection{แอปพลิเคชันต้นแบบมุ่งใช้งานบนระบบปฏิบัติการ Android}
\subsubsection{ประสิทธิภาพในการติดตามการวิ่งอาจแตกต่างกันตามรุ่นและประสิทธิภาพของสมาร์ตโฟน}

% -----------------------------------------------------------
\section{ปัญหาอุปสรรค และแนวทางแก้ไข}
% -----------------------------------------------------------

ปัญหาที่พบระหว่างดำเนินงานและแนวทางแก้ไขแสดงดังตารางที่ \ref{tab:problems}

\begin{table}[H]
\centering
\caption{ปัญหาอุปสรรคและแนวทางแก้ไข}
\label{tab:problems}
\begin{tabular}{|c|p{5.5cm}|p{5.5cm}|}
\hline
\textbf{ลำดับ} & \textbf{ปัญหาที่พบ} & \textbf{แนวทางแก้ไข} \\
\hline
1 & พัฒนาแอปสำหรับติดตามและบันทึกการเดิน–วิ่งของผู้ใช้ & ออกแบบโครงสร้างข้อมูลและแบ่งการทำงานของระบบให้เป็นโมดูล \\ \hline
2 & การกำหนดแผนฝึกมีเงื่อนไขหลายระดับ & กำหนดกฎและเงื่อนไขของแต่ละระดับการฝึกให้ชัดเจน \\ \hline
3 & การพัฒนา UI หลายขั้นตอนอาจทำให้ผู้ใช้สับสน & ออกแบบลำดับการใช้งานให้ต่อเนื่องและลดขั้นตอนที่ไม่จำเป็น \\
\hline
\end{tabular}
\end{table}

% -----------------------------------------------------------
\section{ข้อเสนอแนะ ในการพัฒนาต่อไป}
% -----------------------------------------------------------

\subsection{การปรับปรุงฟังก์ชันการทำงาน}

ควรเพิ่มความแม่นยำในการวิเคราะห์ข้อมูลการวิ่ง ปรับปรุงการแนะนำแผนฝึกให้เหมาะสมกับผู้ใช้แต่ละบุคคล และพัฒนาประสบการณ์ใช้งานให้สะดวกยิ่งขึ้น

\subsection{แนวทางการวิจัยต่อยอด}

ควรทดลองใช้งานกับกลุ่มผู้ใช้จริงในระยะเวลาที่นานขึ้น เพื่อประเมินประสิทธิผลของระบบต่อความสม่ำเสมอและแรงจูงใจในการออกกำลังกาย
